\documentclass[11pt]{article}
\usepackage[a4paper,margin=1in]{geometry}
\usepackage{amsmath,amssymb}
\usepackage[T1]{fontenc}
\usepackage{lmodern}
\usepackage{microtype}
\usepackage{xcolor}
% Footnotes inside \parbox (used below) are otherwise trapped; savenotes emits them
% at the end of the wrapper so they appear on the page foot as usual.
\usepackage{footnote}

\setlength{\parindent}{0pt}
\setlength{\parskip}{0.5em}

\newcommand{\promptbox}[1]{%
  \begin{savenotes}%
  \begingroup\setlength{\fboxsep}{10pt}\noindent
  \colorbox{blue!8}{\parbox{\dimexpr\linewidth-2\fboxsep\relax}{\textbf{Prompt. }#1}}
  \endgroup
  \end{savenotes}%
}
\newcommand{\resultbox}[1]{%
  \begin{savenotes}%
  \begingroup\setlength{\fboxsep}{10pt}\noindent
  \colorbox{purple!8}{\parbox{\dimexpr\linewidth-2\fboxsep\relax}{\textbf{Result. }#1}}
  \endgroup
  \end{savenotes}%
}
\newcommand{\examplebox}[1]{%
  \begin{savenotes}%
  \begingroup\setlength{\fboxsep}{10pt}\noindent
  \colorbox{green!8}{\parbox{\dimexpr\linewidth-2\fboxsep\relax}{\textbf{Example. }#1}}
  \endgroup
  \end{savenotes}%
}

\title{\textbf{Book 1 --- Lecture 2 --- Exercise 3}\\[2mm]
\large Chain Rule for Three Composed Functions}
\author{}
\date{}

\begin{document}
\maketitle

\promptbox{Use the chain rule to find the derivatives of each of the following
functions.\footnote{The \emph{chain rule}: if \(y=h(u)\) and \(u=k(t)\) are
differentiable, then \(\dfrac{dy}{dt}=h'(u)\,k'(t)\). The same idea applies when
the outer function is exponential, logarithmic, or built from sine and cosine.}
\[
\begin{aligned}
g(t) &= \sin(t^2) - \cos(t^2), \\
\theta(\alpha) &= e^{3\alpha} + 3\alpha\,\ln(3\alpha), \\
x(t) &= \sin^2(t^2) - \cos(t^2).
\end{aligned}
\]
}

\section*{\(g(t)=\sin(t^2)-\cos(t^2)\)}
Set \(u(t)=t^2\), so \(\dfrac{du}{dt}=2t\).\footnote{Here the ``inner'' function is
\(t^2\); every appearance of \(t^2\) in \(\sin\) and \(\cos\) is differentiated
through this factor \(2t\).} Then \(g(t)=\sin u-\cos u\), and
\[
\frac{dg}{dt}=\cos u\cdot\frac{du}{dt}-(-\sin u)\cdot\frac{du}{dt}
=(\cos u+\sin u)\frac{du}{dt}.
\]
Substitute \(u=t^2\):
\[
g'(t)=2t\bigl(\cos(t^2)+\sin(t^2)\bigr).
\]

\section*{\(\theta(\alpha)=e^{3\alpha}+3\alpha\,\ln(3\alpha)\)}
Require \(\alpha>0\) so that \(\ln(3\alpha)\) is defined.\footnote{The argument of
\(\ln\) must be positive: here \(3\alpha>0\), equivalently \(\alpha>0\).}

For \(e^{3\alpha}\), let \(v=3\alpha\); then \(\dfrac{d}{d\alpha}e^{3\alpha}
= e^v\cdot 3 = 3e^{3\alpha}\).

For \(3\alpha\,\ln(3\alpha)\), use the product rule with \(a(\alpha)=3\alpha\) and
\(b(\alpha)=\ln(3\alpha)\).\footnote{\(\dfrac{d}{d\alpha}\ln(3\alpha)=\dfrac{1}{3\alpha}\cdot 3
=\dfrac{1}{\alpha}\) by the chain rule on \(\ln(3\alpha)\).}
Now \(a'=3\) and \(b'=1/\alpha\), so
\[
\frac{d}{d\alpha}\bigl(3\alpha\,\ln(3\alpha)\bigr)=3\ln(3\alpha)+3\alpha\cdot\frac{1}{\alpha}
=3\ln(3\alpha)+3.
\]
Therefore
\[
\theta'(\alpha)=3e^{3\alpha}+3\ln(3\alpha)+3
=3\bigl(e^{3\alpha}+\ln(3\alpha)+1\bigr).
\]

\section*{\(x(t)=\sin^2(t^2)-\cos(t^2)\)}
Write \(\sin^2(t^2)=\bigl(\sin(t^2)\bigr)^2\). Again set \(u=t^2\) with
\(\dfrac{du}{dt}=2t\).\footnote{This combines two uses of the chain rule: squaring
the sine, and the argument \(t^2\) inside the sine.} First,
\[
\frac{d}{dt}\sin^2(t^2)=2\sin(t^2)\cos(t^2)\cdot\frac{d}{dt}(t^2)
=4t\,\sin(t^2)\cos(t^2).
\]
Also \(\dfrac{d}{dt}\bigl(-\cos(t^2)\bigr)=\sin(t^2)\cdot 2t=2t\sin(t^2)\). Hence
\[
x'(t)=4t\,\sin(t^2)\cos(t^2)+2t\sin(t^2)
=2t\,\sin(t^2)\bigl(2\cos(t^2)+1\bigr).
\]

\resultbox{Summarizing (with \(\alpha>0\) for \(\theta\)),
\[
\begin{aligned}
g'(t) &= 2t\bigl(\cos(t^2)+\sin(t^2)\bigr),\\
\theta'(\alpha) &= 3e^{3\alpha}+3\ln(3\alpha)+3,\\
x'(t) &= 2t\,\sin(t^2)\bigl(2\cos(t^2)+1\bigr).
\end{aligned}
\]}

\examplebox{At \(t=0\): \(g'(0)=0\) and \(x'(0)=0\) because every term carries a
factor \(t\). At \(\alpha=1\): \(\theta'(1)=3e^3+3\ln 3+3\).}

\end{document}
